\chapter{Résoudre des problèmes}\label{ch:g3:problems}

Connaître les quatre opérations est une chose~; savoir
\emph{laquelle} une histoire demande en est une autre. Ce chapitre
donne une méthode pour les problèmes énoncés --- les mêmes
quatre étapes à chaque fois --- et apprend à lire des
informations dans des tableaux et des pictogrammes.

\section{Les quatre étapes}

\begin{method}[Résoudre un problème énoncé]\label{met:g3:problems:steps}
\begin{enumerate}
  \item \emph{Lire} le problème deux fois~; dire avec ses propres
        mots ce qui est demandé, et ce qui est connu.
  \item \emph{Choisir} l'opération~: dessiner un croquis rapide ou
        une image en barres si le choix n'est pas évident.
  \item \emph{Calculer}, en colonnes ou de tête, avec une estimation
        comme filet de sécurité.
  \item \emph{Répondre} par une phrase complète, avec l'unité
        --- et vérifier que la réponse a du sens (un enfant ne
        pèse pas $300$~kg).
\end{enumerate}
\end{method}

\begin{example}[Quelle opération~?]\label{ex:g3:problems:which}
Les mots donnent un indice sur l'opération, le croquis décide~:
\begin{itemize}
  \item \emph{mettre ensemble / en tout}~: addition~;
  \item \emph{enlever / combien de plus / ce qui manque}~:
        soustraction~;
  \item \emph{autant de groupes de tant / tant de fois}~:
        multiplication~;
  \item \emph{partager équitablement / combien de groupes}~:
        division.
\end{itemize}
\end{example}

\begin{omfigure}
\begin{tikzpicture}[scale=0.62]
  % addition bar
  \draw[thick, fill=omDef!20] (0,0) rectangle (3.4,0.8);
  \draw[thick, fill=omThm!20] (3.4,0) rectangle (5.8,0.8);
  \node[font=\small] at (1.7,0.4) {$26$};
  \node[font=\small] at (4.6,0.4) {$17$};
  \draw[Stealth-Stealth] (0,-0.35) -- (5.8,-0.35);
  \node[below, font=\small] at (2.9,-0.4) {? en tout~: $26 + 17$};
  % comparison bar
  \begin{scope}[xshift=9.4cm]
  \draw[thick, fill=omDef!20] (0,0.9) rectangle (5.2,1.7);
  \node[font=\small] at (2.6,1.3) {$52$};
  \draw[thick, fill=omThm!20] (0,0) rectangle (3.5,0.8);
  \node[font=\small] at (1.75,0.4) {$35$};
  \draw[Stealth-Stealth] (3.5,0.4) -- (5.2,0.4);
  \node[below, font=\small] at (2.6,-0.15)
    {? de plus~: $52 - 35$};
  \end{scope}
\end{tikzpicture}

{\small Les images en barres rendent l'opération visible~: des
parties côte à côte appellent une addition~; deux barres
comparées appellent une soustraction.}
\end{omfigure}

\begin{example}[Un problème en deux étapes]\label{ex:g3:problems:twostep}
\guillemotleft{} Maya achète $3$ paquets de $6$ bouteilles et en
boit $2$. Combien de bouteilles restent-elles~? \guillemotright{}
\begin{enumerate}
  \item Ce qui est demandé~: le nombre de bouteilles restantes.
        Connu~: $3$ paquets de $6$~; $2$ bues.
  \item Croquis~: $3$ groupes de $6$, puis enlever $2$. Deux
        opérations~!
  \item Calcul~: $3 \times 6 = 18$, puis $18 - 2 = 16$.
  \item Phrase~: il reste $16$ bouteilles.
\end{enumerate}
Ne jamais s'arrêter après la première opération~: relire la
question pour voir si l'histoire est finie.
\end{example}

\begin{example}[Un nombre inutile]\label{ex:g3:problems:trap}
\guillemotleft{} Un boulanger, âgé de $52$~ans, cuit $120$ petits
pains et en vend $85$. Combien de petits pains restent-ils~?
\guillemotright{} L'âge est là pour distraire~: la réponse n'a
besoin que de $120 - 85 = 35$ petits pains. Avant de calculer,
toujours trier les nombres~: lesquels la question utilise-t-elle
vraiment~?
\end{example}

\section{Lire des tableaux et des pictogrammes}

\begin{example}[Un tableau]\label{ex:g3:problems:table}
Livres empruntés à la bibliothèque~:
\begin{center}
\renewcommand{\arraystretch}{1.3}
\begin{tabular}{l|cccc}
 & bandes dessinées & romans & sciences & total \\
\hline
semaine 1 & $14$ & $9$ & $6$ & $29$ \\
semaine 2 & $11$ & $12$ & $8$ & $31$ \\
\end{tabular}
\end{center}
Lecture~: semaine 2, romans~: $12$. Calcul à partir du tableau~:
sur les deux semaines ensemble, $14 + 11 = 25$ bandes dessinées ont
été empruntées~; la semaine la plus chargée était la
semaine 2 ($31 > 29$).
\end{example}

\begin{example}[Un pictogramme]\label{ex:g3:problems:pictogram}
Dans un pictogramme, un symbole représente une quantité fixe ---
lire d'abord la légende~!

\begin{omfigure}
\begin{tikzpicture}[scale=0.62]
  \node[left, font=\small] at (0,2.4) {lundi};
  \foreach \i in {0,1,2} \fill[omDef] (0.5+\i*0.8,2.4) circle (0.26);
  \node[left, font=\small] at (0,1.2) {mardi};
  \foreach \i in {0,1,2,3,4} \fill[omDef] (0.5+\i*0.8,1.2) circle (0.26);
  \node[left, font=\small] at (0,0) {mercredi};
  \foreach \i in {0,1} \fill[omDef] (0.5+\i*0.8,0) circle (0.26);
  \fill[omDef] (0.5+1.6,0) circle (0.26);
  \node[right, font=\small] at (5.4,1.2)
    {légende~: un point $= 4$ gâteaux vendus};
\end{tikzpicture}

{\small Un pictogramme~: le lundi montre $3$ points, donc
$3 \times 4 = 12$ gâteaux~; le mardi $5$ points, $20$ gâteaux~;
le mercredi $3$ points, $12$ gâteaux.}
\end{omfigure}
\end{example}

\section{Exercices}

\begin{exercise}[$\star$]\label{exo:g3:problems:1}
Pour chaque histoire, se contenter de \emph{nommer} l'opération
(sans calculer)~: un troupeau de $48$ moutons rejoint un troupeau de
$37$~; \ $56$ cerises partagées entre $7$ enfants~; \ $9$ boîtes
de $8$ crayons~; \ Ali avait $71$ cartes et en a donné $18$.
\end{exercise}

\begin{exercise}[$\star$]\label{exo:g3:problems:2}
Résous avec les quatre étapes~: \guillemotleft{} Un parking a
$246$ voitures le matin~; $158$ de plus arrivent à midi. Combien
de voitures y a-t-il maintenant~? \guillemotright{}
\end{exercise}

\begin{exercise}[$\star$]\label{exo:g3:problems:3}
Résous~: \guillemotleft{} Une corde de $90$~m est coupée en
morceaux de $9$~m. Combien de morceaux~? \guillemotright{}
\end{exercise}

\begin{exercise}[$\star$]\label{exo:g3:problems:4}
Résous~: \guillemotleft{} Les billets coûtent $8$ chacun. Combien
coûtent $26$ billets~? \guillemotright{}
\end{exercise}

\begin{exercise}[$\star$]\label{exo:g3:problems:5}
Dessine une image en barres, puis résous~: \guillemotleft{} Emma a
$64$ autocollants, Hugo en a $29$ de moins. Combien en a Hugo~?
\guillemotright{}
\end{exercise}

\begin{exercise}[$\star$]\label{exo:g3:problems:6}
Deux étapes~: \guillemotleft{} Une caisse contient $45$ pommes. Un
épicer reçoit $6$ caisses et en vend $178$. Combien de pommes
restent-elles~? \guillemotright{}
\end{exercise}

\begin{exercise}[$\star$]\label{exo:g3:problems:7}
Avec le tableau de \cref{ex:g3:problems:table}~: combien de livres de
sciences ont été empruntés sur les deux semaines ensemble~?
Combien de romans de plus en semaine 2 qu'en semaine 1~?
\end{exercise}

\begin{exercise}[$\star$]\label{exo:g3:problems:8}
Avec le pictogramme de \cref{ex:g3:problems:pictogram}~: combien de
gâteaux ont été vendus sur les trois jours ensemble~? Quel jour
en a-t-on vendu le plus~?
\end{exercise}

\begin{exercise}[$\star$]\label{exo:g3:problems:9}
Ce problème a un nombre inutile --- trouve-le, puis résous~:
\guillemotleft{} Un livre de $210$ pages coûte $12$. Léo lit
$35$ pages par jour. Combien de jours lui faut-il pour finir le
livre~? \guillemotright{}
\end{exercise}

\begin{exercise}[$\star\star$]\label{exo:g3:problems:10}
\guillemotleft{} Une famille de $2$ adultes et $3$ enfants visite le
musée. Les billets adultes coûtent $9$, les billets enfants $4$.
Ils paient avec un billet de $50$. Combien de monnaie
reçoivent-ils~? \guillemotright{} (Trois étapes~: adultes,
enfants, puis la monnaie.)
\end{exercise}

\begin{exercise}[$\star\star$]\label{exo:g3:problems:11}
Invente ton propre problème en deux étapes dont la réponse
est $100$, écris sa solution avec les quatre étapes, et teste-le
sur un camarade.
\end{exercise}
